% theme: quadratic_equations_pattern_applications
\MajorQuestionLesson{規則性の利用問題}
\SetRowSpace{6mm}

\TypeQuestionSingle{}{kisei1}
\FigProblemBlock{35mm}{fig/lesson_11_kisei_01}{\Text{碁石を正方形になるように並べ、右上から順に{\rm L}字型に区切っていく。小さいものから順に1番目、2番目、3番目、$\cdots$ とする。}}
{
\QQ{5番目の{\rm L}字型に含まれる碁石の個数を求めなさい。}{}
\QQ{$n$ 番目の{\rm L}字型に含まれる碁石の個数を、$n$ を使って表しなさい。}{}
\QQ{1番目から $n$ 番目までの{\rm L}字型に含まれる碁石の個数の合計を、$n$ を使って表しなさい。}{}
\QQ{1番目から $n$ 番目までの{\rm L}字型に含まれる碁石の個数の合計が64個になるとき、$n$ の値を求めなさい。}{}
}

\TypeQuestionSingle{}{kisei2}
\FigProblemBlock{35mm}{fig/lesson_11_kisei_02}{\Text{正方形のタイルを図のように、1番目は縦1個、横5個、2番目は縦2個、横6個、3番目は縦3個、横7個、$\cdots$ のように、長方形になるよう並べていく。}}
{
\QQ{4番目の長方形に使うタイルの個数を求めなさい。}{}
\QQ{$n$ 番目の長方形に使うタイルの個数を、$n$ を使って表しなさい。}{}
\QQ{$n$ 番目の長方形に使うタイルの個数が77個になるとき、$n$ の値を求めなさい。}{}
\QQ{$n$ 番目の長方形に使うタイルの個数が、外側の一周にあるタイルの個数より60個多くなるとき、$n$ の値を求めなさい。}{}
}

\LessonPageBreak

\TypeQuestionSingle{}{kisei3}
\FigProblemBlock{35mm}{fig/lesson_11_kisei_03}{\Text{図のように、碁石を階段の形に並べる。1段の階段には1個、2段の階段には $1+2$ 個、3段の階段には $1+2+3$ 個、$\cdots$ の碁石を使う。}}
{
\QQ{6段の階段に使う碁石の個数を求めなさい。}{}
\QQ{$n$ 段の階段に使う碁石の個数を、$n$ を使って表しなさい。}{}
\QQ{$n$ 段の階段を2つ作ると、使う碁石の個数は72個になった。$n$ の値を求めなさい。}{}
\QQ{$n$ 段の階段を作ったあと、さらに111個の碁石を加えると、21段の階段に使う碁石の個数と同じになった。$n$ の値を求めなさい。}{}
}

\TypeQuestionSingle{}{kisei4}
\FigProblemBlock{35mm}{fig/lesson_11_kisei_04}{\Text{白いタイルを正方形になるように並べ、そのまわりを黒いタイルで1列だけ囲む。1番目は白いタイルを1辺1個の正方形に並べ、2番目は白いタイルを1辺2個の正方形に並べる。}}
{
\QQ{5番目で使う黒いタイルの個数を求めなさい。}{}
\QQ{$n$ 番目で使う黒いタイルの個数を、$n$ を使って表しなさい。}{}
\QQ{$n$ 番目に使う白いタイルの個数が、$(n+1)$ 番目に使う白いタイルの個数より29枚少ないとき、$n$ の値を求めなさい。}{}
\QQ{$n$ 番目で使う白いタイルの個数が、黒いタイルの個数より56枚多いとき、$n$ の値を求めなさい。}{}
}
